\chapter{Evaluation: measuring quality}

How do we know the model is any good? The codebase computes three families of metrics ---
detection, polygon, and distance --- in \file{eval/}. This chapter is a plain-language
glossary, so the numbers in TensorBoard and \code{tools/eval} mean something.

\section{Detection metrics (COCO)}
These follow the standard COCO protocol (\file{eval/coco\_metrics.py}).

\begin{center}
\renewcommand{\arraystretch}{1.25}
\begin{tabular}{l >{\raggedright\arraybackslash}p{9.4cm}}
\toprule
\textbf{Metric} & \textbf{Meaning} \\
\midrule
\code{mAP} & mean Average Precision averaged over IoU thresholds 0.50:0.95 (the headline COCO number) \\
\code{mAP50} & AP at a single, looser IoU threshold of 0.5 \\
\code{AR100} & average recall with up to 100 detections per image \\
\code{F1score50} & macro-averaged peak F1 at IoU 0.5 --- \textbf{the best-checkpoint selection metric} \\
\code{precision50 / recall50} & precision / recall at each class's peak-F1 operating point \\
\bottomrule
\end{tabular}
\end{center}

\begin{teacher}
\textbf{Precision, recall, F1.} Precision = ``of the boxes I predicted, how many were
right?'' Recall = ``of the real objects, how many did I find?'' They trade off as you move
the confidence threshold. \textbf{F1} is their harmonic mean; the ``peak F1'' is the best
balance over the whole curve. \textbf{Average Precision} summarizes the entire
precision--recall curve as one number; \textbf{mAP} averages it over classes and IoU
thresholds.
\end{teacher}

\begin{why}
The best checkpoint is chosen by \code{F1score50}, not \code{mAP}. F1 at a single operating
point reflects deployable detection quality (a threshold you would actually ship), and it
is more stable run-to-run for this multi-task model than the stricter mAP average. The
per-category table (\code{-{}-per\_category}) exposes which few classes drag the macro
average down.
\end{why}

\section{Polygon metrics}
Predictions are matched to ground truth by \textbf{box} IoU $> 0.5$, then the polygon
\textbf{mask} IoU is computed by rasterizing both outlines (\code{cv2.fillPoly}).

\begin{center}
\renewcommand{\arraystretch}{1.25}
\begin{tabular}{l >{\raggedright\arraybackslash}p{9.4cm}}
\toprule
\code{poly\_mIoU} & mean polygon mask IoU over matched pairs --- pure segmentation quality \\
\code{poly\_recall50} & fraction of GT objects matched at box IoU 0.5 \\
\bottomrule
\end{tabular}
\end{center}

\noindent The polygon conf gate ($0.4$) is a single shared constant
(\code{eval/polygon\_metrics.DEFAULT\_POLY\_CONF\_THRESH}) used by both the metric and the
visualization overlays, so they can never disagree.

\section{Distance metrics}
Each GT is matched to its highest-IoU detection ($\geq 0.5$); errors are in meters
(predictions are $\exp$'d from log-space first), split near/far at 5\,m.

\begin{center}
\renewcommand{\arraystretch}{1.2}
\begin{tabular}{l >{\raggedright\arraybackslash}p{8.6cm}}
\toprule
\code{dist\_mae} & mean absolute error, meters \\
\code{dist\_rmse} & root-mean-squared error (penalizes large misses) \\
\code{dist\_absrel} & mean relative error $|\text{pred}-\text{gt}|/\text{gt}$ \\
\code{dist\_*\_near / *\_far} & the same, restricted to $<5$\,m / $\geq5$\,m \\
\bottomrule
\end{tabular}
\end{center}

\begin{pitfall}
Distance is only scored when the eval split carries distance labels. The shipped distance
dataset is \textbf{training-only}, so during normal validation the distance head is
\emph{not} evaluated --- a distance regression would not show up in val metrics. The
evaluator code exists; it simply isn't wired into the default val loop.
\end{pitfall}

\section{Running it}
\begin{lstlisting}[language=bash]
python -m tools.eval \
    --config configs/experiments/yolo/yolov8_poly_dist.yaml \
    --checkpoint /run/ckpt-100000 --split val --per_category
\end{lstlisting}

\noindent Evaluation automatically uses the EMA weights. \code{tools/continuous\_eval.py}
can watch a run directory and evaluate each new checkpoint automatically, appending to
\code{eval\_log.jsonl}.
